% !TEX program = xelatex
% Bilingual Abstract + Introduction, draft_paper_0414_v4 (BNN rewrite)
% Compile: xelatex intro_abstract_v4.tex
\documentclass[11pt,a4paper]{article}

\usepackage[margin=2.4cm]{geometry}
\usepackage{xeCJK}
\usepackage{amsmath,amssymb}
\usepackage{unicode-math}
\usepackage{setspace}
\usepackage{parskip}
\usepackage{titlesec}
\usepackage{xcolor}

% Fonts — fall back gracefully if specific fonts are unavailable
\setCJKmainfont{PingFang SC}[AutoFakeBold=true]
\setmainfont{Times New Roman}

\titleformat{\section}{\normalfont\large\bfseries}{\thesection}{0.6em}{}
\titleformat{\subsection}{\normalfont\normalsize\bfseries}{\thesubsection}{0.5em}{}

\setlength{\parskip}{0.6em}
\setlength{\parindent}{0pt}
\onehalfspacing

% Convenience macros for the physical symbols the manuscript needs uniformly
\newcommand{\keff}{\ensuremath{k_{\mathrm{eff}}}}
\newcommand{\abase}{\ensuremath{\alpha_{\mathrm{base}}}}
\newcommand{\aslope}{\ensuremath{\alpha_{\mathrm{slope}}}}
\newcommand{\Eint}{\ensuremath{E_{\mathrm{intercept}}}}
\newcommand{\kref}{\ensuremath{k_{\mathrm{ref}}}}
\newcommand{\kslope}{\ensuremath{k_{\mathrm{slope}}}}
\newcommand{\PICP}{\ensuremath{\mathrm{PICP}_{90}}}
\newcommand{\MPIW}{\ensuremath{\mathrm{MPIW}_{90}}}

\title{\textbf{Bayesian Neural Surrogates Unify Uncertainty Propagation,\\
Sensitivity Attribution and Calibration for Coupled\\
Multi-Physics Reactor Analysis}\\[0.4em]
\large 贝叶斯神经代理统一耦合多物理反应堆分析中的\\
不确定性传播、敏感性归因与标定}
\author{Yinuo Chen, Sicheng Wang, Xiaojing Liu, Tengfei Zhang\textsuperscript{*}}
\date{Draft v4 (Round-9 framework-rewrite) \textperiodcentered{} 2026-04-17}

\begin{document}
\maketitle

\section*{Abstract / 摘要}

We train a Bayesian neural network (BNN) surrogate on coupled
OpenMC--FEniCS data of a heat-pipe-cooled reactor and use its posterior
predictive distribution as the single probabilistic object on which
forward uncertainty propagation, Sobol sensitivity attribution and MCMC
posterior calibration all operate, without retraining or recalibration
between them. Two-way coupling damps stress variability by
$\approx 28\%$ relative to a decoupled prediction; stress and
reactivity are governed by largely separate parameters (\Eint{} and
\abase{}); and across $18$ benchmark calibrations the posterior
contracts toward observation-compatible regions with $90\%$-CI coverage
$0.917$. Each posterior predictive draw costs
$\approx 1.76\times 10^{5}\,\times$ less than a coupled high-fidelity
solve ($2266$~s, AMD EPYC 9654), enabling the full uncertainty workflow
to run end-to-end on a single workstation.

本文在耦合 OpenMC--FEniCS 热管冷却反应堆数据上训练贝叶斯神经网络
（BNN）代理，以其后验预测分布作为前向不确定性传播、Sobol 敏感性归因
与 MCMC 后验标定三类分析共用的概率对象，三者之间无需各自重训或重新
校准。双向耦合使应力变异相对解耦预测降低约 $28\%$；应力与反应性分别
由 \Eint{} 与 \abase{} 主导；在 18 例基准标定中，后验分布收敛至与观
测相容的区域，$90\%$-CI 覆盖率为 $0.917$。每次后验预测抽样的代价较
耦合高保真求解（$2266$~s，AMD EPYC 9654）降低约
$1.76\times 10^{5}$ 倍，完整的不确定性工作流可在单台工作站上端到端
运行。

\section{Introduction / 引言}

% --- Paragraph 1: Microreactor background ---
Microreactors --- compact, transportable fission systems rated at
1--20\,MW$_\mathrm{e}$ --- are being developed for decentralized power
generation in remote communities, defense installations and space
missions\textsuperscript{1,11}. NASA and DOE are advancing fission
surface-power systems for sustained lunar and planetary
operations\textsuperscript{12}, while rapidly growing electricity demand
from AI-intensive data centres reinforces the case for distributed,
dispatchable zero-carbon generation\textsuperscript{13}. Unlike
conventional large-scale plants, these designs concentrate fission power,
heat removal and structural loading within a compact monolithic core
where neutronics, heat transfer and structural mechanics are tightly
coupled through two-way feedback\textsuperscript{1,2}. A single
uncertain material parameter can simultaneously alter the neutronic
balance, the thermal field and the stress state; the closed feedback
loop may amplify or damp the resulting variability in ways that
single-physics analyses cannot capture. The uncertainty problem in these
systems is therefore inherently coupled.

微型反应堆——额定 1--20\,MW$_\mathrm{e}$ 的紧凑、可运输裂变系统——正面向
偏远社区、国防设施及太空任务中的分布式供能而开发\textsuperscript{1,11}。
NASA 与 DOE 正推进月面及行星任务的裂变地表供电系统\textsuperscript{12}，
与此同时 AI 密集型数据中心电力需求的快速增长进一步强化了分布式、
可调度零碳发电的需求\textsuperscript{13}。与常规大型电厂不同，
此类设计将裂变功率、热量排出与结构载荷集中于紧凑的一体式堆芯，
中子学、传热与结构力学通过双向反馈紧密耦合\textsuperscript{1,2}。
单个不确定材料参数可同时改变中子平衡、温度场与应力状态；闭合反馈回路
可能以单物理分析无法捕捉的方式放大或衰减由此产生的变异。因此，
此类系统的不确定性问题本质上是耦合的。

% --- Paragraph 2: State of the field + gap ---
High-fidelity coupled solvers now exist for these systems: MOOSE-based
toolchains handle neutronics--thermo-mechanics--heat-pipe coupling for
KRUSTY-class and larger designs\textsuperscript{2,3}, and thermal--structural
analyses have established stress concentration as a persistent safety
concern in monolithic-core configurations\textsuperscript{4,5}.
Surrogate-assisted probabilistic modules in the MOOSE
ecosystem\textsuperscript{6,7} and digital-twin
pipelines\textsuperscript{8} address forward propagation, sensitivity or
calibration individually, but each uses a separately constructed model
with no shared probabilistic basis between analyses.

高保真耦合求解器已适用于这类系统：基于 MOOSE 的工具链可处理 KRUSTY 类
及更大设计上的中子--热力--热管耦合\textsuperscript{2,3}；热--力分析已确认
应力集中是基体式堆芯构型中持续性的安全关注\textsuperscript{4,5}。
MOOSE 生态\textsuperscript{6,7} 与数字孪生管线\textsuperscript{8}
中的代理辅助概率模块分别处理前向传播、敏感性或标定，但每类分析各自
构建独立模型，三者之间缺乏共享的概率基础。

% --- Paragraph 3: Why surrogates needed ---
Running the full probabilistic workflow --- forward uncertainty
propagation, global sensitivity analysis and Bayesian calibration ---
directly on the coupled solver is impractical at ${\sim}10^{3}$\,s per
solve. Gaussian-process and polynomial-chaos surrogates face
dimensionality constraints in nonlinear coupled settings;
physics-informed neural networks require a closed residual form that
iterative solver exchange does not supply; deterministic neural networks
produce only point predictions. None of these provides the coherent
predictive distribution that all three analyses require.

在耦合求解器上直接运行完整的概率工作流——前向不确定性传播、全局
敏感性分析与贝叶斯标定——不可行，每次求解约需 $10^{3}$\,s。高斯过程
与多项式混沌代理在非线性耦合设置下受维度限制；物理信息神经网络需要
闭合残差形式，而迭代求解器交换无法提供；确定性神经网络仅产生点估计。
上述方案均无法提供三类分析共同所需的连贯预测分布。

% --- Paragraph 4: BNN solution ---
A Bayesian neural network (BNN) fills this gap. By placing a
distribution over the network weights, a BNN yields a posterior
predictive distribution that can be sampled at negligible cost once
trained. This single distribution serves as the shared probabilistic
layer on which forward propagation, variance-based sensitivity
attribution and observation-conditioned posterior calibration all
operate, with no retraining or recalibration between them.

贝叶斯神经网络（BNN）填补了这一空缺。通过对权重引入分布，BNN 在
训练完成后即可以极低代价反复采样其后验预测分布。该分布充当前向传播、
基于方差的敏感性归因与观测条件后验标定共同使用的概率层，三者之间无需
重训或重新校准。

% --- Paragraph 5: Framework declaration ---
This paper presents a BNN-based posterior-predictive framework for
end-to-end probabilistic analysis of coupled multi-physics systems. The
framework unifies forward uncertainty propagation, Sobol sensitivity
decomposition and MCMC posterior calibration on a single predictive
layer trained on coupled OpenMC--FEniCS data of a heat-pipe-cooled
reactor. It reveals that two-way coupling damps stress variability by
$\approx 28\%$, that stress and reactivity are governed by largely
separate parameter pathways, and that the posterior calibration contracts
the material-parameter distribution toward observation-compatible
regions --- all at five orders of magnitude less cost per evaluation
than a coupled high-fidelity solve.

本文提出了一套基于 BNN 后验预测分布的耦合多物理系统端到端概率分析
框架。该框架将前向不确定性传播、Sobol 敏感性分解与 MCMC 后验标定统一
到基于热管冷却反应堆耦合 OpenMC--FEniCS 数据训练的单一预测层上。框架
揭示：双向耦合使应力变异降低约 $28\%$；应力与反应性分别受制于基本独立
的参数通道；后验标定将材料参数分布收敛至与观测相容的区域——以上分析
的单次评估代价较耦合高保真求解低约五个数量级。

\end{document}
